% =====================================================================
% Appendix E -- Conveyors from the literature (reduced table).
% The table is generated by maps/paper_tables.py (table_cases_reduced.tex);
% the full table, with provenance and notes, is Table S1 of the
% supplementary material (table_cases.tex).
% =====================================================================
\section{Conveyors from the literature}\label{app:cases}

\Cref{tab:cases} lists the published conveyors used in Section~\ref{sec:results}, with the
quantities that support Section~\ref{sec:betaregime}. Each value is computed with the code of the
paper from the data of its source. The mass ratio $\beta$ follows from the design take-up tension
through Eq.~\eqref{eq:beta_tension} where the source gives it (S, SM, Pa, Lo, Si), since it can be
checked against the static balance; otherwise from the published counterweight mass (H, LL, G,
WR) or, for Su, from the mass that holds the published slack-side tension. Taking the
counterweight as hung directly on the $\nstr$ strands ($\lambda=1$) unless the rigging is
published gives an upper bound on $\beta$ for a roped rig. Where both the take-up tension and the
counterweight mass are published (Si), the rigging follows from $\nstr\Tt=\lambda M_wg$:
$\lambda=0.5$ with $\nstr=2$, equivalent to a counterweight hung directly on four strands. For S
the published take-up tension cannot be held by the published counterweight; the table keeps the
tension, which gives the larger $\beta$ (\cref{app:validation} and the note of Table~S1). The full table of the supplementary material
(Table~S1) adds the densities, wave speeds, belt speeds, take-up tensions, rigging, the criterion
that set the take-up tension according to each source and the provenance of every input.

% Generated by maps/paper_tables.py from maps/cases.py; do not edit by hand.
\begin{table}[tbp]
\centering
\footnotesize
\setlength{\tabcolsep}{4pt}
\begin{threeparttable}
\caption{Conveyors from the literature, by decreasing length. $\xi$: take-up position from the drive exit along the return strand, in units of $L$; $\gamma=c_r/c_c$ at full material coupling and, in brackets, at the lowest coupling of the material class \citep{lodewijks2002}; $t_1$: fundamental period with the take-up mass; $W_r=g\mu_rL$; $\beta=(4/n)\lambda T_t/W_r$; $\beta_5$: take-up mass that lengthens the fundamental by 5\,\% at the case's $\xi$ and $\gamma$.}
\label{tab:cases}
\begin{tabular}{@{}lrclrlll@{}}
\toprule
Case & $L$ (km) & $\xi$ & $\gamma$ & $t_1$ (s) & $T_t/W_r$ & $\beta$ & $\beta/\beta_5$ \\
\midrule
Si & 13.10 & 0.001 & 2.24$^{\mathrm{c}}$ (1.84) & 66.1$^{\mathrm{c}}$ & 0.025$^{\mathrm{c}}$ & 0.025$^{\mathrm{c}}$ & 0.03$^{\mathrm{c}}$ \\
LL & 7.60 & 0.005 & 1.97 (1.72) & 46.8$^{\mathrm{a}}$ & 0.042$^{\mathrm{a}}$ & 0.084$^{\mathrm{a}}$ & 0.11$^{\mathrm{a}}$ \\
S & 7.12 & 0.001 / 0.999 & 1.67 (1.47) & 40.4 / 29.0 & 0.066 & 0.131 & 0.22 / 0.27 \\
H & 5.10 & 0.005 & 0.97$^{\mathrm{m}}$ & 28.4 & 0.012 & 0.012 & 0.03 \\
G & 4.50 & 0.001 & 2.20 (1.77) & 24.6$^{\mathrm{a}}$ & -- & 0.006$^{\mathrm{a}}$ & 0.01$^{\mathrm{a}}$ \\
SM & 3.00 & 0.100 & 1.95 (1.28) & 13.6$^{\mathrm{a}}$ & 0.065 & 0.131$^{\mathrm{a}}$ & 0.17$^{\mathrm{a}}$ \\
Pa & 2.56 & 0.020 & 1.85 (1.30) & 9.5 & 0.084 & 0.169 & 0.24 \\
NC & 2.48 & 0.037 & 2.46 & 27.5 & -- & -- & -- \\
Lo & 1.00 & 0.001 & 2.76 (2.15) & 28.5 & 0.102 & 0.205 & 0.15 \\
Su & 0.81 & 0.010 & 2.58 (2.01) & 6.1$^{\mathrm{c}}$ & 0.318 & 0.635 & 0.42 \\
WR & 0.27 & -- & 2.93$^{\mathrm{c}}$ (2.21) & 5.2--5.9$^{\mathrm{a}}$ & 0.442$^{\mathrm{a}}$ & 0.884$^{\mathrm{a}}$ & 0.56--0.58$^{\mathrm{a}}$ \\
\bottomrule
\end{tabular}
\begin{tablenotes}[flushleft]
\footnotesize
\item Superscripts: weakest provenance among the inputs of each value (c, belt or idler catalogue \citep{continental2021,fennerdunlop2009,cema}; a, assumed; m, measured); unmarked values are published or derived from published data. Densities, wave speeds, take-up tensions, rigging, the criterion that set $T_t$ and the notes on each case are in Table~S1 of the supplementary material. S uses the take-up tension published by its source, which its 4500~kg counterweight could not hold; with that counterweight $\beta=0.017$ and the take-up mass lengthens $t_1$ by 0.15\,\% at the tail.
\item Si, \citet{sinaga2008kpc}; LL, \citet{li2009amesim}; S, \citet{song2012}; H, \citet{harrison1983,harrison1985b}; G, \citet{gao2026}; SM, \citet{suchorab2025longdistance}; Pa, \citet{pascual2005}; NC, \citet{nordell1984}; Lo, \citet{lodewijks1996}; Su, \citet{surtees1995runback}; WR, \citet{wheatley2021}.
\end{tablenotes}
\end{threeparttable}
\end{table}

